\section{Problem Definition}\label{sec:problem}

\subsection{Witnesses of non-equivalence}
Let $\mathcal S$ be a relational schema given as SQLite DDL with declared constraints $\Sigma$: primary keys, \texttt{NOT NULL}, \texttt{UNIQUE} and foreign keys. For a query string $Q$ and a database $D$, $Q(D)$ is the list of rows SQLite returns, or an error. A \emph{comparison contract} $\mathcal C$ decides whether two results are equal. Our primary contract is $\mathcal C_{\mathrm{set}}$, defined as $\mathrm{set}(R_1)=\mathrm{set}(R_2)$, which is the test under which BIRD defines execution accuracy. We also report multiset equality $\mathcal C_{\mathrm{bag}}$ and a lenient contract $\mathcal C_{\mathrm{len}}$, which equates an empty result with one all-\texttt{NULL} row and compares numbers to nine significant digits.

Given a set $\mathcal A$ of admissible databases, $Q_1$ and $Q_2$ are \emph{equivalent under $(\mathcal C,\mathcal A)$} if $\mathcal C(Q_1(D),Q_2(D))$ holds for every $D\in\mathcal A$. A database $W\in\mathcal A$ with $\neg\mathcal C(Q_1(W),Q_2(W))$ is a \emph{witness} of non-equivalence. Membership in $\mathcal A$ and the contract are both decided by executing SQLite on $W$. A witness is therefore a certificate that anyone can replay. Its absence, however, proves nothing.

Given the benchmark database $D_0$, a prediction $Q_p$ is \emph{EX-correct} if $\mathcal C_{\mathrm{set}}(Q_p(D_0),\allowbreak Q_g(D_0))$ for the gold query $Q_g$. Conversely, every EX-incorrect prediction that executes is non-equivalent, since $D_0$ separates it. These \emph{known positives} measure the recall of a disprover without manual labels.

\subsection{Admissibility tiers}
A database that satisfies $\Sigma$ can still be absurd, e.g., a \texttt{DATE} column holding \texttt{\textquotesingle{}abc\textquotesingle{}}. We use the following nested tiers:
\begin{itemize}
\item $\mathcal A_0$: $D$ satisfies $\Sigma$ and both queries execute.
\item $\mathcal A_1\subseteq\mathcal A_0$: type-faithful values, non-\texttt{NULL} primary keys and non-empty referenced tables. Type-faithful values are numbers in numeric columns and valid ISO-8601 text in declared date columns.
\item $\mathcal A_{1f}\subseteq\mathcal A_1$: \emph{format-faithful} databases. Text values follow a lexical shape observed for their column, in which digits keep their count and letter runs collapse (e.g., \texttt{9999-99-99 99:99:99}). Temporal text columns hold ISO-like text.
\item $\mathcal A_2\subseteq\mathcal A_{1f}$: referenced columns satisfy the range, categorical, not-null and dependency constraints mined from the official databases and validated with an LLM by SpotIt+~\cite{tremante2026spotitplus}.
\end{itemize}
$\mathcal A_0$--$\mathcal A_2$ are derived from the schema and published constraint sets. They are the tiers under which existing verifiers report their counts.

\subsection{Data-profile admissibility}\label{sec:profile}
Benchmark maintainers have the benchmark database $D_0$, and its data constrain which databases are plausible far more than its DDL does. We define tier $\mathcal A_R\subseteq\mathcal A_{1f}$ by a \emph{data profile} $\Pi$ that standard profiling queries extract from $D_0$~\cite{abedjan2015profiling}:
\begin{itemize}
\item \emph{Column domains}: whether \texttt{NULL} occurs, the numeric range and integrality, and the observed values of categorical columns. A text column is categorical if it has $\le$100 distinct values or $\le$1\,\% of rows, and a numeric column if it has $\le$20 distinct values. For other columns, the domain records lexical shapes, letter case and the absence of surrounding blanks. Key columns may hold new identifiers of the observed type.
\item \emph{Unique column combinations} of one or two identifier-like columns. Reals, dates and many-valued measures are excluded, as their uniqueness is accidental.
\item \emph{Inclusion dependencies}: undeclared foreign keys between equally named columns, total participation (every parent row has a child row), and symmetric and irreflexive self-relations.
\item \emph{Lookup tables}, whose contents are fixed. A lookup table has $\le$200 rows and $\le$3 columns and is referenced by a foreign key.
\item \emph{Value dependencies}: functional dependencies between categorical columns, one-to-one maps between identifier columns, cross-key dependencies (a child column that equals a parent column on $\ge$99.5\,\% of joined rows), and year--date agreement.
\end{itemize}
A witness $W$ is in $\mathcal A_R$ if every value lies in its column's domain, \texttt{NULL} occurs only in columns that contain \texttt{NULL} in $D_0$, and $W$ satisfies the dependencies and lookup contents of $\Pi$. $\mathcal A_{Rq}$ checks these conditions only on columns referenced by either query. On BIRD, mining $\Pi$ for eleven databases with 3.9\,M rows takes 258\,s on one core (Table~\ref{tab:profile}). The 66\,s for value dependencies were timed in a re-run that reproduced the frozen profile exactly.

\subsection{Meaningful witnesses}\label{sec:meaningful}
A witness shows that two queries differ on some admissible database, not that the prediction is wrong. With respect to the question, the difference can expose a \emph{prediction error}, a \emph{gold error} or an \emph{ambiguity} of the question. Alternatively, it can arise only on \emph{implausible data}. Such data satisfies every recorded constraint but does not occur in the benchmark instance and, in the annotators' judgment, would not occur in the domain the question is about (e.g., a played match without a score). The last judgment goes beyond the instance and is made by annotators with read access to $D_0$ (Section~\ref{sec:attribution}). We call a witness \emph{meaningful} if it falls into one of the first three classes. The share of implausible witnesses measures how far a tier is from the data the questions are about. We call an EX-correct prediction with a witness a \emph{certified EX false positive}: it is certified non-equivalent. As Section~\ref{sec:attrresults} shows, this is not the same as being wrong.
